\mode*
\section{The staging area}
\label{sec:staging}

The second educational objective is the three places a change can be:
working tree, staging area (index), repository---and the commands that move
changes between them.

\subsection{Documented misconceptions and difficulties}

The staging area is a prime candidate for a concept misfit: the analysis in
\textcite{DeRossoJackson2016Gitless} redesigns Git largely around
difficulties of this kind, and course observations report students
struggling with what \texttt{add} accomplishes
\autocite{IsomottonenCochez2014}.
More specifically, \textcite{IsomottonenCochez2014} document \texttt{add}
misread as \emph{make}: learners take \texttt{git add file} to be
\textcquote[Sec.~5.3]{IsomottonenCochez2014}{associated with creating a new
file, while it stages the file for the commit}, so that \enquote{one creates
a new file and then commits it} rather than moving an existing change from
the working tree into the index.
The same authors note a habit of \texttt{git add .}\ and \texttt{git commit
-a} that keeps working tree and index in lock-step, hiding the very
distinction the staging area exists to make \autocite{IsomottonenCochez2014}.

\subsection{Candidate critical aspects}
\label{sec:staging-aspects}

The staging area is where the documented difficulties and the design
analysis agree most, so two of the three candidates rest on the
literature; the third is posited.

\begin{description}
  \item[Three places, two moves]
    Working tree, index, history; \texttt{add} moves a change from the
    first to the second, \texttt{commit} from the second to the third.
    Read off the documented reading of \texttt{add} as \emph{make}
    \autocite[Sec.~5.3]{IsomottonenCochez2014}---a learner who has not
    discerned the index has no second place for \texttt{add} to move
    anything into, so the command must be doing something else---and
    off the hidden dependency that an unadded file is unknown to Git
    \autocite{Church2014HiddenDependencies}.
    Tested by the item \emph{what \texttt{git add} does}
    (\cref{app:knowledge-items}); the open item \emph{keeping one of two
    changes} (\cref{app:open-items}) sets the task the index exists for.
  \item[The same file can differ in all three places]
    A file can be staged in one state and edited to another at once.
    Posited from the object of learning; the \texttt{add .}/\texttt{commit
    -a} habit that keeps tree and index in lock-step
    \autocite{IsomottonenCochez2014} is what a learner does who has not
    discerned it, and the dependency between an edit and the index is
    among those professionals find hard to hold
    \autocite{Church2014HiddenDependencies}.
    Tested by the item \emph{two versions at once}.
  \item[\texttt{git status} reads the three places]
    Status reports a file by where it differs, not by whether it is
    saved.
    Posited from the object of learning: it is the instrument through
    which the other two aspects are observed at all.
    Tested by the item \emph{what status shows after an edit}.
\end{description}

\subsection{Patterns of variation}

\begin{itemize}
  \item Contrast: modify a file; observe \texttt{git status} and
    \texttt{git diff}/\texttt{git diff --staged} after editing, after
    \texttt{add}, after a second edit---the same file in three states.
  \item Generalisation: the add--commit cycle across different kinds of
    change (new file, modification, deletion).
  \item Fusion: stage part of the changes and commit while the working
    tree still differs.
\end{itemize}
